\section{实现语言与框架}
\label{sec:framework}

我们完全使用嵌入Python的CuTe-DSL~\citep{cutedsl}编写\faf，没有任何CUDA C++组件。
CuTe-DSL编译器接受Python源代码，降级到PTX，然后使用PTX编译器（ptxas）最终生成汇编代码（SASS）。

\textbf{完整表达能力与简洁抽象。}
CuTe-DSL编程模型与CUTLASS C++同构，确保\faf在受益于Python元编程生产力提升和快速JIT编译时间的同时，保留了底层GPU编程的完整表达能力。CuTe-DSL提供对PTX的直接访问作为逃生舱，允许开发者实现任何所需功能而不受框架限制。例如，我们利用自定义PTX序列实现CuTe-DSL API中尚未完整暴露的操作（这些将在未来版本中集成），证明了我们的框架不会将开发者限制在有限的GPU能力子集中。

\textbf{通过JIT实现快速编译。}
由于复杂的C++模板元程序，编译时间一直是过去FlashAttention实现的瓶颈。通过将CuTe-DSL嵌入Python并采用即时（JIT）编译，\faf比传统基于C++模板的方法实现了更快的构建速度。如表~\ref{table:compile_time}所示，\faf比\fat将编译时间缩短了20--30倍。这种快速迭代周期显著提升了开发者的生产力，在内核开发过程中支持更快的实验和调试。

\begin{table}[h]
  \centering
  \begin{tabular}{lcc}
    \toprule
    方法 & 前向传播 & 后向传播 \\
    \midrule
    \fat & 55s & 45s \\
    \faf & 2.5s & 1.4s \\
    \midrule
    加速比 & 22$\times$ & 32$\times$ \\
    \bottomrule
  \end{tabular}
  \caption{单个内核的编译时间：FA3（C++模板）与FA4（CuTe-DSL）。通常FA2和FA3需要为不同注意力变体预编译数百个内核。}
  \label{table:compile_time}
\end{table}

\textbf{灵活性与可访问性。}
基于Python的框架在实践中已展现出其灵活性：开发者已成功在\faf之上构建FlexAttention和块稀疏注意力变体，而无需修改核心框架。通过降低入门门槛，我们的方法使具有几个月GPU编程经验的研究人员和工程师在无需C++模板元编程深厚专业知识的情况下，即可贡献有意义的扩展。这种可访问性加速了创新，使注意力机制研究社区能够更快速地探索新的算法变体。

我们的愿景是提供一个综合框架，以最优性能构建各种注意力变体。\faf不是从头实现每个注意力变体，而是将公共功能分解为独立的可组合原语。块稀疏模式、掩码策略、可变序列长度处理和工作调度等操作均作为正交原语暴露，可以自由组合。这种模块化设计确保优化和新特性惠及所有基于框架构建的注意力实现，同时通过编译为高效GPU内核仍然实现最高性能。
